\chapter{MATLAB 与 MEX 架构及代码映射}
\label{chap:code}

\section{总体架构}
当前仓库采用“MATLAB 调度 + C/MEX 核心算子”的混合架构。MATLAB 负责一切需要可读性和灵活性的逻辑，例如配置读取、粒子初始化、主循环、restart、日志和图形输出；C/MEX 则负责粒子对上的高频循环，把最耗时的邻居搜索和离散算子下沉到编译层。与早期版本相比，当前的 MATLAB 侧已经收拢为单个主入口 \code{SPH\_Poiseuille.m}，并在文件内部通过 \code{local functions} 组织流程，同时把运行态显式封装为 \code{cfg}、\code{state}、\code{geom}、\code{neighbor} 和 \code{monitor} 五类结构。物理 MEX 使用源文件 \code{mex/sph\_physics\_mex.c}，运行时输出名为 \code{sph\_physics\_shell\_mex}；它除了保留底层物理 mode 外，还新增了 \code{advance\_shell\_step} 与 \code{wall\_shear\_monitor} 两个高层门面接口。

\begin{table}[H]
    \centering
    \caption{主要文件的职责划分}
    \label{tab:file-responsibility}
    \begin{tabular}{p{0.28\textwidth} p{0.18\textwidth} p{0.42\textwidth}}
        \toprule
        文件 & 主要位置 & 职责 \\
        \midrule
        \code{SPH\_Poiseuille.m} & 行 1--757 & 单主文件 MATLAB 壳层：配置解析、结构体组织、主循环、restart、后处理和绘图 \\
        \code{build\_shell\_wall\_particles.m} & 行 1--20 & 生成上下单层 shell 壁面及其几何元数据 \\
        \code{mex/sph\_neighbor\_search\_mex.c} & 行 1--337 & cell-linked list 邻居搜索、周期最小像距、核函数与导数计算 \\
        \code{mex/sph\_physics\_mex.c} & 行 1--1456 & 底层物理 mode + 高层门面 \code{advance\_shell\_step}/\code{wall\_shear\_monitor} \\
        \code{config.ini} & 行 1--19 & 物理参数与仿真控制参数配置 \\
        \bottomrule
    \end{tabular}
\end{table}

\section{主脚本的阶段映射}
主脚本的实现边界相当清晰，可以直接按阶段阅读：
\begin{enumerate}[leftmargin=2.5em]
    \item 行 16--36：初始化路径、处理输出覆盖环境变量并自动编译两个 MEX 文件；
    \item 行 38--84：读取配置并计算派生物理量；
    \item 行 86--117：初始化流体粒子、单层 shell 壁面粒子和基础状态数组；
    \item 行 119--155：做 restart 签名检查和状态恢复；
    \item 行 157--234：做初始邻居搜索、密度修正，并将运行数据收拢为 \code{cfg/state/geom/neighbor/monitor}；
    \item 行 236--294：执行统一时间步主循环；
    \item 行 296--508：做后处理、误差评估和图形输出；
    \item 行 510--757：定义局部辅助函数，例如环境变量覆盖、邻居缓存、MEX 门面桥接、时间步、排序、周期包裹与壁面钳制。
\end{enumerate}
这种布局对报告写作非常友好，因为每个阶段几乎都能映射成一小节，不需要在多个文件之间来回跳转；同时，结构体组织也把旧版本里散落的大量并列变量压缩到了更稳定的数据边界上。

\section{物理 MEX 的模式分发}
\code{mex/sph\_physics\_mex.c} 的设计方式仍然是“一个 MEX，多个 mode”。MATLAB 调用时首先传入一个字符串模式名，MEX 再把请求分发给具体函数。当前支持七个模式：
\begin{enumerate}[leftmargin=2.5em]
    \item \code{density\_correction}：密度重初始化、体积计算与核梯度修正矩阵；
    \item \code{viscous\_force}：流体-流体和流体-壁面的层流粘性力；
    \item \code{transport\_correction}：粒子位置 shifting；
    \item \code{integration\_1st}：第一阶段积分；
    \item \code{integration\_2nd}：第二阶段积分；
    \item \code{advance\_shell\_step}：复用底层 mode 的单步推进门面；
    \item \code{wall\_shear\_monitor}：上下壁面剪应力统计门面。
\end{enumerate}
这种模式分发结构的好处在于，MATLAB 侧只需维护一个物理 MEX 句柄，不必在主脚本中散落多个二进制接口；同时，不同算子之间又能共享粒子对、体积和修正矩阵的数据组织方式。与旧版“MATLAB 直接拼接五个底层 mode”相比，当前主循环实际只直接依赖 \code{advance\_shell\_step} 和 \code{wall\_shear\_monitor} 两个高层接口，底层 mode 主要留作内部复用和验证。

\section{代码与公式的一一对应}
表 \ref{tab:formula-map} 列出了报告中最关键的“公式—代码”对应关系。它不是完整索引，但足以支撑定位主线逻辑。

\begin{table}[H]
    \centering
    \caption{核心公式与代码段对应}
    \label{tab:formula-map}
    \begin{tabular}{p{0.22\textwidth} p{0.28\textwidth} p{0.38\textwidth}}
        \toprule
        模块 & 关键公式 & 代码位置 \\
        \midrule
        驱动力计算 & 式 \eqref{eq:body-force} & \code{SPH\_Poiseuille.m} 行 65--73 \\
        统一时间步 & 式 \eqref{eq:dt} & \code{SPH\_Poiseuille.m} 行 637--643 \\
        shell 壁面几何 & 式 \eqref{eq:single-shell-wall} & \code{build\_shell\_wall\_particles.m} 行 10--19 \\
        密度重初始化 & 式 \eqref{eq:density-correction} & \code{mex/sph\_physics\_mex.c} 行 89--362 \\
        核梯度修正 & $B_i=w_1P_i+w_2I$ & \code{mex/sph\_physics\_mex.c} 行 89--362 \\
        粘性项 & 式 \eqref{eq:viscous-wall} 及其流体-流体版本 & \code{mex/sph\_physics\_mex.c} 行 384--538 \\
        transport correction & 式 \eqref{eq:transport-correction} & \code{mex/sph\_physics\_mex.c} 行 556--696 \\
        第一阶段积分 & 式 \eqref{eq:rho-half-stage1}--\eqref{eq:half-position} & \code{mex/sph\_physics\_mex.c} 行 717--947 \\
        第二阶段积分 & 式 \eqref{eq:position-full-finish}--\eqref{eq:rho-half-finish} & \code{mex/sph\_physics\_mex.c} 行 964--1097 \\
        单步推进门面 & 五个底层算子的组合调度 & \code{mex/sph\_physics\_mex.c} 行 1116--1326 \\
        壁面剪应力监测 & 上下壁面 pair-wise 剪应力统计 & \code{mex/sph\_physics\_mex.c} 行 1339--1429 \\
        防穿透 & 距壁钳制到 $0.5dp$ & \code{SPH\_Poiseuille.m} 行 739--757 \\
        \bottomrule
    \end{tabular}
\end{table}

\section{Restart、排序与输出模块}
\code{SPH\_Poiseuille.m} 除了数值积分外，还承担了三个重要的工程模块。

其一是 restart。程序在每个输出点保存完整的粒子状态、时间和步数。恢复时会检查配置签名是否匹配，而当前签名字符串已经显式包含 \code{wall=single-shell-noslip}，因此不会再出现旧版本那种“边界模式切换但签名没变”的问题。与此同时，主脚本还支持通过环境变量覆盖 \code{config.ini} 与输出路径，这使得 smoke test 和隔离目录批处理可以不污染默认结果文件。

其二是排序优化。函数 \code{sort\_particles\_by\_cell} 会按照局部网格编号重新排列流体粒子，从而提升后续 MEX 访问时的缓存局部性。排序后壁面粒子保持在数组尾部，这样流体粒子数量 \code{n\_fluid} 的语义始终清晰。

其三是输出。主脚本不只生成最终速度场图，还持续记录中截面速度剖面的时间演化，并保存 restart 文件。对数值方法报告来说，这部分信息很宝贵，因为它让读者看到的不只是“终态对不对”，还有“收敛过程是不是合理、运行能否断点续算”。当前实现中，主循环已经不再显式串接五个底层算子，而是以 \code{advance\_one\_step\_mex} 和 \code{wall\_shear\_monitor\_mex} 为桥接点，把高层流程控制和底层粒子对计算分开。

\section{文档与实现的偏差}
在本次同步后，README 与主报告的主要口径已经与源码对齐；当前更容易再次过时的，不再是名称本身，而是报告中写死的代码行号、mode 数量和主循环阶段描述。只要后续继续重构 \code{SPH\_Poiseuille.m} 或 \code{mex/sph\_physics\_mex.c}，这一章和附录里的索引就需要同步刷新。另一个仍然存在的工程偏差是：\code{mex/sph\_neighbor\_search\_mex.c} 的文件头写有“支持 OpenMP 并行”，但源码主体并没有实际的 \code{pragma omp} 并行循环。上述问题不影响程序当前运行，却会直接影响后续论文或技术文档的准确性，因此本报告统一以源码真实行为为准。
